


\subsection{Structural Model Analysis: Estimation with Three Types}
\label{sec: appendix3types}
We re-estimated the model by introducing a third student type to assess whether accounting for additional time-invariant unobserved heterogeneity substantially alters the simulated outcomes. The results suggest this is not the case.
As shown in Figure \ref{fig:types_fraction}, the distribution of student types remains largely unchanged when moving from a two-type to a three-type specification. In the latter model, only 4 percent of students are assigned to the third type.

The model continues to have a good fit in terms of students' choices and outcomes (Table \ref{tab:descriptionoutcomes3types}), as well as for most main treatment effects (Tables \ref{TEprecollegeoutcomes3types} and \ref{fitTEadmenrper3types}).



\begin{figure}[H]
	\centering
\begin{subfigure}[b]{0.49\textwidth}
\includegraphics[width=\linewidth]{output/figures/type_histogram_2types.png}
\caption{Two types} 
\end{subfigure}
\begin{subfigure}[b]{0.49\textwidth}
\includegraphics[width=\linewidth]{output/figures/type_histogram_3types.png}
\caption{Three types} 
\end{subfigure}
\caption{\label{fig:types_fraction} This figure shows the fraction of students of type 1 and 2 estimated in a model with two types (Panel A) and the fraction of students of type 1, 2 and 3 estimated in a model with three types (Panel B).} 
\end{figure}


\input{output/tables/summary_stats_outcomes_3types} 


\input{output/tables/TE_pre_college_outcomes_3types} 



\input{output/tables/fit_TE_moments_adm_enr_per_3types}

